% =====================================================================
%  PCB DESIGN ENGINEER PORTFOLIO
%  Vladyslav Lysiuk
%
%  Англійська версія portfolio.tex — дослівний переклад, без нового змісту.
%  Структура, макроси, зображення й висоти ілюстрацій ті самі; змінено лише
%  мову тексту та \setdefaultlanguage.
%
%  Як зібрати:  xelatex portfolio_EN.tex   (двічі — заради змісту/полів)
%
%  Рушій ОБОВ'ЯЗКОВО XeLaTeX або LuaLaTeX. pdflatex НЕ підійде.
% =====================================================================

\documentclass[11pt,a4paper]{article}

% ---------- Поля сторінки (15 мм з усіх боків) ----------
\usepackage[a4paper,margin=15mm,bottom=18mm,footskip=8mm]{geometry}

% ---------- Шрифти та мова ----------
\usepackage{fontspec}
\setmainfont{DejaVu Sans Condensed}
\newfontfamily\headingfont{DejaVu Sans Condensed}

\usepackage{polyglossia}
\setdefaultlanguage{english}

% ---------- Графіка, кольори, допоміжні пакети ----------
\usepackage{graphicx}
\usepackage{xcolor}
\usepackage[most]{tcolorbox}
\usepackage{enumitem}          % компактні списки
\usepackage{microtype}
\usepackage{fancyhdr}
\usepackage{ragged2e}

% ---------- ПАЛІТРА (та сама, що й в українській версії) ----------
\definecolor{accent}{HTML}{0F5C4A}   % акцент — темно-зелений
\definecolor{ink}{HTML}{14181D}      % основний текст
\definecolor{muted}{HTML}{5D6772}    % сірий для метаданих і підписів
\definecolor{chipbg}{HTML}{EEF3F1}   % заливка «чипів»
\definecolor{line}{HTML}{D7DDE3}     % тонкі лінії та рамки

\color{ink}

% ---------- Колонтитул: номер сторінки внизу праворуч ----------
\pagestyle{fancy}
\fancyhf{}
\renewcommand{\headrulewidth}{0pt}
\fancyfoot[R]{\footnotesize\color{muted}\thepage}

\setlength{\parindent}{0pt}
\setlength{\parskip}{4pt}

\widowpenalty=10000
\clubpenalty=10000

% =====================================================================
%  ВЛАСНІ КОМАНДИ
% =====================================================================

\newtcbox{\chip}{on line, boxsep=0pt, left=5pt, right=5pt, top=3pt, bottom=3pt,
  colframe=line, colback=chipbg, boxrule=0.4pt, arc=3pt,
  fontupper=\footnotesize\color{accent!85!black}}

% Шапка складається з трьох окремих абзаців у спільній групі: інакше заборона
% переносів не діє (TeX читає \hyphenpenalty аж наприкінці абзацу).
\newcommand{\projecthead}[3]{%
  \begingroup
  \setlength{\parskip}{0pt}%
  \hyphenpenalty=10000\exhyphenpenalty=10000\relax
  \raggedright
  {\footnotesize\color{muted}\addfontfeature{LetterSpace=25}\MakeUppercase{#1}\par}
  \vspace{3pt}
  {\LARGE\bfseries\color{ink}#2\par}
  \vspace{4pt}
  {\footnotesize\color{muted}#3\par}
  \endgroup
  \vspace{4pt}
  \hrulecolor{line}{0.6pt}
  \vspace{8pt}
}

\newcommand{\chips}[1]{\par\noindent{\raggedright #1\par}\vspace{2pt}}
\newcommand{\cs}{\hskip3pt\penalty0\hskip0pt\relax}
\newcommand{\hrulecolor}[2]{\par\noindent{\color{#1}\rule{\linewidth}{#2}}\par}

\newtcolorbox{rolebox}{enhanced, breakable=false, colback=white,
  boxrule=0pt, frame hidden, sharp corners,
  borderline west={2pt}{0pt}{accent},
  left=10pt, right=0pt, top=5pt, bottom=5pt, boxsep=0pt}

\newcommand{\role}[1]{%
  \par\vspace{6pt}
  \begin{rolebox}
    \small\justifying{\bfseries\color{accent}My role.} #1
  \end{rolebox}
  \par\vspace{4pt}
}

\newcommand{\tz}{\item}
\newcommand{\notes}[1]{%
  \par\vspace{2pt}
  {\footnotesize\color{muted}%
   \begin{itemize}[leftmargin=12pt, itemsep=1pt, topsep=2pt, parsep=0pt,
     label=\textcolor{muted}{\rule{3pt}{3pt}}]
     #1
   \end{itemize}}%
  \par\vspace{2pt}
}

\newenvironment{points}
  {\begin{itemize}[leftmargin=14pt, itemsep=3pt, topsep=4pt, parsep=0pt,
     label=\textcolor{accent}{\small$\bullet$}]\small\justifying}
  {\end{itemize}}

\newcommand{\bigimg}[3][0.42]{%
  \par\vspace{6pt}
  \noindent
  \begin{minipage}{\linewidth}
    \centering
    \setlength{\fboxrule}{0.4pt}\setlength{\fboxsep}{2pt}%
    \fcolorbox{line}{white}{\includegraphics[width=\dimexpr\linewidth-6pt\relax,
                                            height=#1\textheight,
                                            keepaspectratio]{#2}}\\[3pt]
    {\footnotesize\color{muted}#3}
  \end{minipage}
  \par\vspace{8pt}
}

\newcommand{\pairimg}[5][0.30]{%
  \par\vspace{6pt}
  \noindent
  \begin{minipage}{\linewidth}
    \setlength{\fboxrule}{0.4pt}\setlength{\fboxsep}{2pt}%
    \begin{minipage}[t]{0.485\linewidth}
      \centering
      \fcolorbox{line}{white}{\includegraphics[width=\dimexpr\linewidth-6pt\relax,
                                              height=#1\textheight,
                                              keepaspectratio]{#2}}\\[3pt]
      {\footnotesize\color{muted}#3}
    \end{minipage}\hfill
    \begin{minipage}[t]{0.485\linewidth}
      \centering
      \fcolorbox{line}{white}{\includegraphics[width=\dimexpr\linewidth-6pt\relax,
                                              height=#1\textheight,
                                              keepaspectratio]{#4}}\\[3pt]
      {\footnotesize\color{muted}#5}
    \end{minipage}
  \end{minipage}
  \par\vspace{8pt}
}

\usepackage{calc}

\graphicspath{{images/}}

% =====================================================================
%  ДОКУМЕНТ
% =====================================================================
\begin{document}

% Документ починається одразу з першого проєкту — без шапки з контактами
% та без вступного абзацу. Проєкти впорядковані хронологічно.

% ---------------------------------------------------------------
%  PROJECT 01 — ATtiny13 thermometer (course project)
% ---------------------------------------------------------------
\projecthead{Project 01}%
  {Thermometer on ATtiny13}%
  {Course project in “Modern CAD for electronic equipment” $\cdot$ 2 layers $\cdot$ 2023}

\chips{\chip{Altium Designer}\cs\chip{2 layers}\cs\chip{ATtiny13}\cs\chip{First board}}

\notes{
  \tz The first board I did on my own in Altium Designer.
  \tz Purpose - a room thermometer: LM35 sensor, 74HC595 shift registers,
      three 7-segment displays.
}

\role{Component library, schematic, board routing, full set of design documentation.}

\bigimg[0.28]{temp_iso.jpg}{3D view in Altium Designer.}
\pairimg[0.30]{temp_3d_top.jpg}{Top view.}%
              {temp_layout.jpg}{Routing, two layers.}

% ---------------------------------------------------------------
%  PROJECT 02 — remote power electronics laboratory board
% ---------------------------------------------------------------
\clearpage
\projecthead{Project 02}%
  {Remote power electronics laboratory board}%
  {Horizon Europe $\cdot$ NGI Search, grant No. 101069364 $\cdot$ 2 layers $\cdot$ batch of 5 $\cdot$ 2023--2024}

\chips{\chip{Altium Designer}\cs\chip{Buck converter}\cs\chip{Horizon Europe}}

\notes{
  \tz The first production project I took part in - the board went the whole way
      from schematic to fabrication and assembly.
}

\role{Component library and schematics in Altium Designer, placement and routing of
the board, full assembly and soldering of all five boards.}


\bigimg[0.44]{buck_3d.jpg}{3D view in Altium Designer.}
\bigimg[0.40]{buck_layout.jpg}{Routing.}
\bigimg[0.40]{ngi_assembled.jpg}{Assembled board.}
\pairimg[0.64]{ngi_batch.jpg}{Fabricated batch.}%
              {ngi_stand.jpg}{Inside the bench together with Raspberry Pi and Analog Discovery.}

\vspace{4pt}
{\footnotesize\color{muted}Carried out within the NGI Search project, grant agreement No. 101069364.}
\vspace{3pt}
\hrulecolor{line}{0.6pt}

% ---------------------------------------------------------------
%  PROJECT 03 — shield for NUCLEO-H723
% ---------------------------------------------------------------
\clearpage
\projecthead{Project 03}%
  {Expansion board for STM32 NUCLEO-H723}%
  {Remote microprocessor technology laboratory $\cdot$ 2 layers, high density $\cdot$ Erasmus+ DIGITRANS $\cdot$ 2024--2025}

\chips{\chip{Altium Designer}\cs\chip{2 layers, high density}\cs\chip{STM32H723}\cs%
\chip{C for MCU}}

\notes{
  \tz Routed on two layers only, in line with the project's manufacturing budget.
  \tz Footprints and board outline matched to the NUCLEO-H723 development bench.
  \tz Five sets assembled by hand (THT + SMD) and integrated into teaching benches
      together with Raspberry Pi and Analog Discovery.
}

\role{Component library, schematic work and board routing done under the guidance of a
lead engineer. Assembly and soldering of all five sets, test firmware.}



\pairimg[0.38]{shield_iso.jpg}{3D view in Altium Designer.}%
              {shield_layout.jpg}{Routing, both layers.}
\pairimg[0.38]{shield_assembled.jpg}{Board after SMD assembly.}%
              {shield_nucleo.jpg}{Mounted on NUCLEO-H723.}
\bigimg[0.48]{shield_stand.jpg}{The bench in operation: Analog Discovery 2, the shield with Nucleo and a Raspberry Pi on a shared panel.}

% ---------------------------------------------------------------
%  PROJECT 04 — thermocouple temperature measurement
% ---------------------------------------------------------------
\clearpage
\projecthead{Project 04}%
  {24-channel thermocouple temperature measurement system}%
  {Industrial measurement board $\cdot$ 4 layers $\cdot$ STM32F4 $\cdot$ fabricated in Germany $\cdot$ 2025--2026}

\chips{\chip{Altium Designer}\cs\chip{STM32}\cs\chip{Galvanic isolation}\cs%
\chip{J/K thermocouples}}

\notes{
  \tz Not a from-scratch development - adaptation of an existing industrial project
      to the customer's new requirements.
}

\role{Rework of the thermocouple measurement path schematic, the channel switching circuit, removal of unused blocks of the project. Changes to the component library and to the printed circuit board.
Done during the internship at Hochschule Bonn-Rhein-Sieg (DAAD Eastern Partnership).}



\bigimg[0.44]{tc_iso.jpg}{3D view in Altium Designer.}
\bigimg[0.40]{tc_top.jpg}{Top view: functional zones.}
\bigimg[0.38]{tc_bottom.jpg}{Bottom view.}

% ---------------------------------------------------------------
%  PROJECT 05 — solenoid valve driver
% ---------------------------------------------------------------
\clearpage
\projecthead{Project 05}%
  {8-channel solenoid valve driver}%
  {Power electronics $\cdot$ 4 layers $\cdot$ STM32 $\cdot$ fabricated, batch of 5 $\cdot$ 2026}

\chips{\chip{Altium Designer}\cs\chip{STM32}\cs\chip{Power electronics}}

\notes{
  \tz Partial circuit calculation and development of the mechanical design and the board
      under the guidance of a lead engineer.
  \tz The placement follows the requirement to separate the power and control sections:
      power circuits around the perimeter, the microcontroller isolated in the centre.
}

\role{Partial calculation of component values, component library, full routing of the printed circuit board. Work under a lead engineer's supervision. The basis of the bachelor's qualification work.}


\bigimg[0.46]{valve_3d_iso.jpg}{3D view in Altium Designer.}
\bigimg[0.38]{valve_3d_top.jpg}{Top view: mirror symmetry of the eight power channels, the control section in the centre.}
\bigimg[0.46]{valve_board.jpg}{The fabricated board (batch of 5) before component assembly.}

% ---------------------------------------------------------------
%  PROJECT 06 — Arduino Duemilanove compatible board
% ---------------------------------------------------------------
\clearpage
\projecthead{Project 06}%
  {Arduino Duemilanove compatible board}%
  {Study project $\cdot$ 2 layers $\cdot$ ATmega328 $\cdot$ THT + SMD}

\chips{\chip{Altium Designer}\cs\chip{2 layers}\cs\chip{ATmega328}\cs%
\chip{USB}\cs\chip{5\,V supply}\cs\chip{Mixed assembly}}

\notes{
  \tz Study project: recreating a board compatible with Arduino Duemilanove -
      ATmega328 in a DIP package, USB bridge.
  \tz Mixed assembly chosen for hand soldering: SMD on the top layer,
      the microcontroller and connectors - THT.
}

\role{Building the component library and routing the board, a study project.}


\pairimg[0.32]{ard_iso.jpg}{3D view in Altium Designer.}%
              {ard_top.jpg}{Top view.}
\bigimg[0.32]{ard_layout.jpg}{Routing, both layers.}

% ---------------------------------------------------------------
%  PROJECT 07 — SiC MOSFET gate driver
% ---------------------------------------------------------------
\clearpage
\projecthead{Project 07}%
  {SiC MOSFET gate driver for a dynamic wireless power transfer system}%
  {Module for a DWPT inverter $\cdot$ 2 layers $\cdot$ 2026}

\chips{\chip{Altium Designer}\cs\chip{SiC MOSFET}\cs\chip{Galvanic isolation}\cs%
\chip{ANSYS}}

\notes{
  \tz The schematic is not my own: the Wolfspeed CGD15SG00D2 reference driver
      for 3rd generation SiC MOSFETs (C3M) was taken as a basis.
  \tz The board design was reworked for the requirements of the inverter of a hybrid
      dynamic wireless power transfer system - the shape and mounting dimensions
      were adapted for modular replacement.
  \tz The driver is built as a separate module soldered into the inverter board: this way
      the unit can be replaced without touching the main board.
  \tz The output stage routing was optimised to reduce parasitic capacitance
      and inductance during fast switching.
}

\role{Adaptation and complete re-layout of the printed circuit board, component library,
routing of the output stage, simulation of parasitic parameters in ANSYS.}

\bigimg[0.44]{sic_iso.jpg}{3D view in Altium Designer.}
\pairimg[0.34]{sic_top.jpg}{Front view.}%
              {sic_bottom.jpg}{Bottom side: digital isolator and gate circuits.}
\bigimg[0.50]{sic_layout.jpg}{Routing: the board outline is fitted to its seat in the inverter board.}

\hrulecolor{line}{0.6pt}

\end{document}
